
\scanpage{58}
\mspage{14}
C) \uline{Découpage} \deleted{\unsure{de}} suivant une partie \inserted{ouv. et f.}
$\Sigma'\subset\Sigma_0$ de la sous-variété formée des cercles de découpage --- cela correspond à une partie $A'$ de $A_0$ (arêtes \deleted{\unsure{ordinaires}} \inserted{liées}), mais cette fois ci il ne s'agit pas de ``gommer'' i.e. de contracter en un point, mais de ``découper'', ce qui au niveau du graphe revient à découper une arête par son milieu (i.e. à la remplacer par deux arêtes libres).

Le \uline{marquage suit}. \deleted{\unsure{Les marq., par un}}
Mais il y a ici à distinguer entre le cas d'une opération anodine \inserted{en $C$} --- celui
\deleted{où il y a \unsure{tel et un}}
$C$ est effectivement marqué i.e. $C\cap R\ne\varnothing$ i.e. $C\subset\Sigma^+$ --- qui à peu de choses près (une ambiguïté dans un ens. fini, \deleted{)} en plus) peut se récupérer, à l'aide du produit de cette opération --- et celui où $C$ n'est pas marqué.
Mais il semble que dans la façon dont nous allons travailler
\mspage{57}
